%=========== ΛΥΣΗ - 42ο ΘΕΜΑ Α ===========
\providecommand{\theoriaref}[3]{%
\begin{center}
\fcolorbox{maincolor!70!black}{maincolor!8}{%
\begin{minipage}{.88\linewidth}
\textcolor{maincolor!80!black}{\kerkissans{\textbf{#1~\ref{#2}}}}\\[1mm]
#3
\end{minipage}}
\end{center}}
\providecommand{\orismosref}[1]{%
\theoriaref{Ορισμός}{#1}{Η απάντηση δίνεται από τον αντίστοιχο ορισμό.}}
\providecommand{\apodeixiref}[1]{%
\theoriaref{Θεώρημα}{#1}{Η ζητούμενη απόδειξη δίνεται στην αντίστοιχη απόδειξη.}}

\begin{thema}{Α}
\begin{erwthma}
\item \apodeixiref{thewrhma:paragwgos_efaptomenhs}

\item \orismosref{orismos:praxeis_synarthsewn}

\item \orismosref{orismos:koilh_synarthsh}

\item Οι χαρακτηρισμοί των προτάσεων είναι:
\begin{alist}
\item \textbf{Σωστό}. Από το θεμελιώδες θεώρημα του ολοκληρωτικού λογισμού, αν $F$ είναι παράγουσα της $f$, τότε
\[ \dint_a^\beta f(x)\d x=F(\beta)-F(a). \]
\item \textbf{Σωστό}. Είναι ο ορισμός της πραγματικής συνάρτησης με πεδίο ορισμού το $A$.
\item \textbf{Λάθος}. Η $\ef x$ δεν ορίζεται στα σημεία όπου $\syn x=0$, δηλαδή στα $x=\dfrac{\pi}{2}+k\pi$, $k\in\mathbb{Z}$.
\item \textbf{Σωστό}. Αν ο στιγμιαίος ρυθμός μεταβολής είναι $0$, τότε τη συγκεκριμένη στιγμή το μέγεθος δεν μεταβάλλεται.
\item \textbf{Σωστό}. Αν η παράγωγος είναι θετική στα διαστήματα ανάμεσα στα μεμονωμένα μηδενικά της, τότε η συνάρτηση παραμένει γνησίως αύξουσα.
\end{alist}

\item Σωστή απάντηση είναι η \textbf{γ}. Από τις σχέσεις
\[ \lim_{x\to+\infty}\dfrac{f(x)}{x}=\lambda
\quad\text{και}\quad
\lim_{x\to+\infty}(f(x)-\lambda x)=\mu \]
προκύπτει ότι η $C_f$ έχει πλάγια ασύμπτωτη
\[ y=\lambda x+\mu \]
στο $+\infty$.
\end{erwthma}
\end{thema}
%=========== ΤΕΛΟΣ ΛΥΣΗΣ - 42ο ΘΕΜΑ Α ===========
